\worksheettitle
  {Точки на отрезке}
  {\modulemark{M04} Учебный модуль \textbullet\ полный маршрут}

\studentnameblock

\begin{notebox}[Чему вы научитесь]
  Определять порядок точек на отрезке, выбирать целое и его части, строить
  схему и находить неизвестную длину.
\end{notebox}

\begin{checkpointbox}[Вход: проверьте опору]
  Точка $C$ лежит между $A$ и $B$. $AB=11$ см, $AC=4$ см.
  \[
    CB=\answerblank[22mm]-\answerblank[22mm]=\answerblank[28mm].
  \]
  Какой отрезок здесь целый? \answerblank[25mm]
\end{checkpointbox}

\begin{notebox}[Как устроен маршрут]
  Сначала прочитайте порядок точек. Затем разберите равенство целого и частей,
  потренируйтесь на одной и двух внутренних точках, исправьте ошибку и
  выполните контроль без готовой схемы.
\end{notebox}

\section{Порядок, целое и части}

\betweenfigure

\begin{fillbox}[Разберите пример]
  На рисунке порядок точек $A-C-B$: точка $C$ лежит между $A$ и $B$.
  Поэтому отрезок $AB$ --- целый, а $AC$ и $CB$ --- его части:
  \[
    AB=AC+CB.
  \]
  Если $AB=13$ см и $AC=5$ см, то
  \[
    CB=AB-AC=13-5=\answerblank[18mm]\text{ см}.
  \]
\end{fillbox}

\begin{notebox}[Алгоритм]
  \textbf{1.} Запишите порядок точек.
  \textbf{2.} Назовите целый отрезок.
  \textbf{3.} Запишите его как сумму последовательных частей.
  \textbf{4.} Выберите действие и вычислите.
  \textbf{5.} Проверьте, что часть не больше целого.
\end{notebox}

\section{Практика с опорой}

\begin{practicebox}[1. Прочитайте порядок]
  \begin{enumerate}[label=\alph*)]
    \item $M-K-N$: между $M$ и $N$ лежит \answerblank[15mm];
    \item $P-T-S$: целый отрезок \answerblank[18mm], его части
      \answerblank[18mm] и \answerblank[18mm].
  \end{enumerate}
\end{practicebox}

\wholepartsfigure

\begin{practicebox}[2. Найдите целое и часть]
  \begin{enumerate}[label=\alph*)]
    \item $AC=3$ см $8$ мм, $CB=5$ см $7$ мм. Найдите $AB$.
    \item $AB=14$ см, $CB=6$ см $5$ мм. Найдите $AC$.
  \end{enumerate}
  \writinglines{3}
\end{practicebox}

\section{Две точки внутри отрезка}

\fourpointsfigure

\begin{fillbox}[Что сделали]
  При порядке $A-C-D-B$:
  \[
    AB=AC+CD+DB,\qquad CD=AB-AC-DB.
  \]
  Если $AB=12$ см, $AC=4$ см, $DB=3$ см, то
  \[
    CD=12-4-3=\answerblank[18mm]\text{ см}.
  \]
\end{fillbox}

\begin{practicebox}[3. Расстояния от одного конца]
  Точки $P$ и $Q$ лежат на $MN$. $MP=3$ см $5$ мм,
  $MQ=8$ см $2$ мм.

  Ближе к $M$ находится \answerblank[15mm].
  \[
    PQ=\answerblank[25mm]-\answerblank[25mm]
      =\answerblank[30mm].
  \]
\end{practicebox}

\section{Самостоятельная практика}

\begin{practicebox}[4. Сначала запишите порядок]
  На отрезке $KL$ отмечены точки $M$ и $N$ в порядке $K-M-N-L$.
  $KL=18$ см, $KM=6$ см, $NL=5$ см.

  Равенство частей: \answerblank[100mm]

  Найдите $MN$. \writinglines{2}
\end{practicebox}

\section{Разбираем ошибку}

\begin{warningbox}[Почему нельзя сразу складывать?]
  Из одной точки $A$ отложены длины $AC=4$ см и $AD=9$ см.
  Ученик написал: $CD=AC+AD=13$ см.

  \textbf{Причина ошибки:} \answerblank[100mm]

  \textbf{Исправление:} \answerblank[100mm]
\end{warningbox}

\section{Контрольная точка}

\begin{checkpointbox}[5. Выполните без подсказки]
  На отрезке $AB$ точки $C$ и $D$ расположены в порядке $A-C-D-B$.
  $AB=20$ см, $AC=7$ см, $DB=4$ см.
  \begin{enumerate}[label=\textbf{\arabic*.}]
    \item Нарисуйте схему. \writinglines{2}
    \item Запишите равенство целого и частей. \answerblank[95mm]
    \item Найдите $CD$. \writinglines{2}
    \item Объясните, как проверили результат. \writinglines{1}
  \end{enumerate}
\end{checkpointbox}

\begin{reflectionbox}[Проверяю себя]
  \begin{itemize}[label=\choicebox]
    \item определяю порядок точек до вычисления;
    \item различаю целый отрезок и его части;
    \item записываю равенство по схеме;
    \item проверяю найденную длину.
  \end{itemize}
\end{reflectionbox}

\begin{resultbox}[Дальнейший маршрут]
  Если путаете порядок или целое и части, выполните M04-R. Для тренировки
  используйте M04-H. При устойчивом результате продолжайте следующий раздел.
  M04\(\star\) выполняется дополнительно.
\end{resultbox}
